\section{Uniform momentum capacity and uncut gradient moments}
\label{sec:v100-capacity}
The preceding pressure theorem applies to the exact optical weight
through the already proved marked comparison. No unweighted word
count or unspecified finite fiber multiplicity enters the conclusion.

\begin{theorem}[Uniform weighted momentum capacity]
\label{thm:v100-capacity}
For the original radius family and every finite collision count,
\begin{equation}\label{eq:v100-capacity}
 \mathcal K_m\le K_*:=
 1+C\left(Q_*+D Q_*^2\frac{2q_0}{1-q_0}\right)<\infty.
\end{equation}
Consequently, for every fixed admissible auxiliary band,
\begin{equation}\label{eq:v100-SC}
       m^6\mathcal K_m\theta_{B_*}^{2m}\longrightarrow0.
\end{equation}
The strict finite-block criterion of Theorem~\ref{thm:v99-block}
is also satisfied analytically for a sufficiently large fixed block
at each such band.
\end{theorem}
\begin{proof}
Insert $P_j(R)\le Q_*$ in Theorem~\ref{thm:v99-pressure-comparison}.
The complete depth sum is at most $2q_0/(1-q_0)$; its coefficient
$D$ counts the single omitted physical step. Global fixed-itinerary
momentum uniqueness and the optical/stable comparison have already
been proved before this insertion. Thus this bounds the actual
$I_m\omega/\det K$ sum on each momentum fiber, including all
internal incidence scales and all exact labels.

Since $\theta_{B_*}<1$, the right side of
$m^6\mathcal K_m\theta_{B_*}^{2m}\le K_*m^6\theta_{B_*}^{2m}$
tends to zero with $B_*$ fixed. Also
$\sup_R P_L(R)\le Q_*$, so that
\[
 \max\{1,(\sup_R P_L(R))^{1/L}\}\theta_{B_*}^2<1
 \quad\text{if}\quad
 L>\frac{\log\max(1,Q_*)}{-2\log\theta_{B_*}}.
\]
This is an existence certificate with finite constants, not a
numerical interval computation of an optimal block or mixing rate.
The block may depend on the already fixed auxiliary band, never on
$m$. Equation~\eqref{eq:v100-SC} itself does not require a block
choice at all.
\end{proof}

Put $\alpha=129/128$. For $1\le p<\alpha$ and $0<b<b_0$, define
an actual nonnegative measure on the same exact event by
\begin{equation}\label{eq:v100-moment-source}
 w_{p,b}(t)\,dt
   =F_\#\!\left(I_m z^{-p}\one_{\{0<z<b\}}\bar\nu\right).
\end{equation}
Here the exponent acts on $z$, not on $I_m$; no assertion about
$(I_m/z)^p$ is made. Values on the original area-null critical set
are immaterial. The source is untruncated at $z=0$.

\begin{lemma}[The two-cap moment integral]\label{lem:v100-layer}
If $E,H,b>0$ and $1/2<p<2$, then
\begin{align}
 &b^{-p}\min\{E b^{1/2},H b^2\}
   +p\int_0^b v^{-p-1}\min\{E v^{1/2},H v^2\}\,dv\notag\\
 &\qquad\le C_p\min\{H b^{2-p},
                 E^{(4-2p)/3}H^{(2p-1)/3}\}.
                                                   \label{eq:v100-layer}
\end{align}
\end{lemma}
\begin{proof}
The quadratic cap alone gives $2H b^{2-p}/(2-p)$.
The caps agree at $v_0=(E/H)^{2/3}$. Below $v_0$ use $Hv^2$;
above it use $Ev^{1/2}$. The two resulting integrals are bounded by
multiples of $Hv_0^{2-p}$ and $Ev_0^{1/2-p}$, respectively, and
both equal $E^{(4-2p)/3}H^{(2p-1)/3}$. The boundary term has the
same bound. Finiteness at zero uses $p<2$, and at the extended
upper limit it uses $p>1/2$. Both inequalities hold together,
proving their minimum after increasing $C_p$.
\end{proof}

\begin{theorem}[Uniform uncut first and nearby gradient moments]
\label{thm:v100-uncut}
For every $1\le p<\alpha$, every fixed auxiliary band $B_*$ and
$|J|=L$, the source in \eqref{eq:v100-moment-source} satisfies
\begin{align}
 m^2\int_J w_{p,b}\le{}&C_p b^{\alpha-p}(L+B_*^{-1})\notag\\
 &+C_p\min\left\{m^2K_*b^{2-p},
 [C_{B_*}(1+B_*L)]^{(4-2p)/3}
 m^2K_*^{(2p-1)/3}\theta_{B_*}^{m(4-2p)/3}\right\}.
                                               \label{eq:v100-uncut-finite}
\end{align}
Constants are uniform in radius, exact labels, collision count and
interval location. In particular, for each fixed $b,L>0$,
\begin{equation}\label{eq:v100-uncut-limit}
 \limsup_m\sup_{R,n,k}\sup_{|J|=L}m^2\int_J w_{p,b}
                                      \le C_p Lb^{\alpha-p}.
\end{equation}
The full source $v^0=F_\#(I_m/z\,\bar\nu)/dt$ obeys
\begin{equation}\label{eq:v100-v0}
 \limsup_m\sup_{R,n,k}\sup_{|J|=L}m^2\int_Jv^0\le CL.
\end{equation}
No spectral-capacity hypothesis remains in these conclusions.
\end{theorem}
\begin{proof}
The endpoint-momentum area formula and
Theorem~\ref{thm:v100-capacity} give
$\int_Ju_v\le K_*v^2$, at every count and every $v>0$ in range.
The inherited critical-volume estimate gives
\[
 m^2\int_Ju_v\le C v^\alpha(L+B_*^{-1})+Ev^{1/2},
 \qquad E=C_{B_*}(1+B_*L)m^2\theta_{B_*}^m.
\]
The source on the left is nonnegative. Combining its two upper
bounds therefore retains the leading term and caps the error by
$\min(Ev^{1/2},Hv^2)$, with $H=m^2K_*$.

Apply the positive identity
\[
 z^{-p}\one_{\{z<b\}}=
 b^{-p}\one_{\{z<b\}}+
 p\int_0^b v^{-p-1}\one_{\{z<v\}}\,dv.
\]
Integrate first from a positive lower endpoint and then use monotone
convergence. The leading integral is finite because $p<\alpha$;
Lemma~\ref{lem:v100-layer} pays the entire error down to zero.
The exponents of $m^2$ in the two factors sum to one, giving
\eqref{eq:v100-uncut-finite}. The physical quadratic bound also
proves finiteness of the measure at each count. Local coarea away
from its area-null critical set, followed by monotone exhaustion,
gives its density.

Fix $b,L,B_*$ before the collision limit. The second expression in
the minimum tends to zero uniformly in $R,n,k$. Enlarge $B_*$ only
afterwards to obtain \eqref{eq:v100-uncut-limit}. For $p=1$, add
the region $z\ge b_0/2$ using the inherited complete $I_m$ source
bound. This proves \eqref{eq:v100-v0}. Neither a collision-dependent
cutoff nor a band growing with $m$ is used.
\end{proof}
